% A translation of Li & Wang, "Collaborative Neurodynamic Algorithms for Solving Sudoku Puzzles"
% (ICIST 2022, pp. 8-17) to n-Queens Completion, instantiated to what is actually proved in
% HopfieldNet/QUBO/Instances/Queens/.
%
% The point of this document is fidelity in two directions at once: it follows the structure and
% register of the Sudoku paper, and it says explicitly wherever the queens case does NOT line up
% with the Sudoku case.  Section 3.4 is the table of those differences; nothing elsewhere claims a
% correspondence that section does not license.
%
% Every board drawn below, every count and every table entry is output of the development:
%   lake build HopfieldNet.QUBO.Instances.Queens.Bench   -- the two tables, ~16 min
%   HopfieldNet/QUBO/Instances/Queens/Bench.lean:80      -- the corpus
% The chessboards use the same normalised outline and palette as the interactive widget, so a figure
% here and a frame in the infoview are the same picture.
%
% NOTE: copy `easychair.cls` and `lstlean.tex` from the main paper's directory next to this file.
\documentclass{easychair}

\usepackage{amsmath}
\usepackage{amssymb}
\usepackage{booktabs}
\usepackage{caption}
\usepackage{listings}
\usepackage{xcolor}
\usepackage{csquotes}
\usepackage{fontawesome}
\usepackage{tikz}

\newtheorem{thm}{Theorem}[section]
\newtheorem{theorem}[thm]{Theorem}
\newtheorem{lemma}[thm]{Lemma}
\newtheorem{definition}{Definition}[section]

\definecolor{academicpurple}{rgb}{0.4,0.0,0.35}
\newcommand{\secref}[1]{\hyperref[#1]{\textcolor{academicpurple}{\S\ref*{#1}}}}
\newcommand{\citeref}[1]{\textcolor{academicpurple}{\cite{#1}}}
\newcommand{\thmref}[1]{\textcolor{academicpurple}{\hyperref[#1]{Theorem~\ref*{#1}}}}
\newcommand{\tabref}[1]{\textcolor{academicpurple}{\hyperref[#1]{Table~\ref*{#1}}}}
\newcommand{\figref}[1]{\textcolor{academicpurple}{\hyperref[#1]{Figure~\ref*{#1}}}}
\newcommand{\extlink}{~\ensuremath{{}^\text{\faExternalLink}}}
% TODO before submission: replace COMMIT with the pushed hash.
\newcommand{\repo}{https://github.com/mkaratarakis/HNBM/blob/COMMIT}
\newcommand{\src}[1]{\href{\repo/#1}{\extlink}}

\definecolor{keywordcolor}{rgb}{0.7, 0.1, 0.1}
\definecolor{tacticcolor}{rgb}{0.0, 0.1, 0.6}
\definecolor{commentcolor}{rgb}{0.4, 0.4, 0.4}
\definecolor{symbolcolor}{rgb}{0.0, 0.1, 0.6}
\definecolor{sortcolor}{rgb}{0.1, 0.5, 0.1}
\definecolor{attributecolor}{rgb}{0.7, 0.1, 0.1}
\definecolor{codebg}{rgb}{0.96,0.96,0.99}
\definecolor{codeframe}{rgb}{0.75,0.70,0.80}
\def\lstlanguagefiles{lstlean.tex}
\lstset{language=lean,
  basicstyle=\fontsize{8.4}{9.8}\selectfont\ttfamily,
  backgroundcolor=\color{codebg},
  frame=single, framerule=0.5pt, rulecolor=\color{codeframe},
  framesep=4pt, xleftmargin=6pt, xrightmargin=6pt,
  aboveskip=4pt, belowskip=4pt}

%%% ------------------------------------------------------------------- chessboards
%%% Palette and queen outline taken from queensBoard.js, so that the figures below and the widget of
%%% Section 6.4 are the same picture.  A square is addressed (column, -row) with row 0 at the top,
%%% matching the representation used throughout: q_i is the column of the queen in row i.
\definecolor{sqlight}{HTML}{F0D9B5}
\definecolor{sqdark}{HTML}{B58863}
\definecolor{sqedge}{HTML}{8A6242}
\definecolor{attackc}{HTML}{C0392B}
\definecolor{qdarkbody}{HTML}{26221E}
\definecolor{qdarkedge}{HTML}{0D0B0A}
\definecolor{qlitebody}{HTML}{FBFAF7}
\definecolor{qliteedge}{HTML}{4A423A}
\definecolor{cRow}{HTML}{1F5FA8}
\definecolor{cCol}{HTML}{1E7A48}
\definecolor{cDiag}{HTML}{B26B00}
\definecolor{cAnti}{HTML}{6B3FA0}

%%% Every macro below is prefixed `qc' so that nothing here can collide with the document class or
%%% with a package loaded later.  `qcpic' opens a picture whose baseline is its centre, so boards of
%%% different sizes line up when set side by side in a tabular.
\newenvironment{qcpic}[1]
  {\begin{tikzpicture}[x=#1,y=#1,baseline=(current bounding box.center)]}
  {\end{tikzpicture}}

% \qcboard{N} -- the squares and the frame of an N x N board.
\newcommand{\qcboard}[1]{%
  \pgfmathtruncatemacro{\qcn}{#1-1}%
  \foreach \qci in {0,...,\qcn}{\foreach \qcj in {0,...,\qcn}{%
    \pgfmathtruncatemacro{\qcmix}{mod(\qci+\qcj,2)==0 ? 100 : 0}%
    \fill[sqlight!\qcmix!sqdark] (\qcj,-\qci) rectangle ++(1,-1);}}%
  \draw[sqedge,line width=0.5pt] (0,0) rectangle (#1,-#1);}

% \qcglyph{col}{row}{body}{edge} -- one queen, from the widget's outline.
\newcommand{\qcglyph}[4]{%
  \begin{scope}[shift={(#1,-#2)},line join=round,line width=0.35pt,draw=#4,fill=#3]
    \filldraw (0.14,-0.255) -- (0.265,-0.45) -- (0.32,-0.185) -- (0.41,-0.45)
      -- (0.50,-0.155) -- (0.59,-0.45) -- (0.68,-0.185) -- (0.735,-0.45)
      -- (0.86,-0.255) -- (0.80,-0.555) -- (0.20,-0.555) -- cycle;
    \filldraw[rounded corners=0.5pt] (0.175,-0.555) rectangle (0.825,-0.640);
    \filldraw (0.255,-0.64) .. controls (0.255,-0.78) and (0.205,-0.83) .. (0.165,-0.875)
      -- (0.835,-0.875) .. controls (0.795,-0.83) and (0.745,-0.78) .. (0.745,-0.64) -- cycle;
    \filldraw[rounded corners=0.5pt] (0.125,-0.875) rectangle (0.875,-0.950);
    \foreach \qcu/\qcv in {0.14/0.235, 0.32/0.165, 0.50/0.135, 0.68/0.165, 0.86/0.235}
      \filldraw (\qcu,-\qcv) circle (0.072);
  \end{scope}}
% A given queen is dark, a queen the solver placed is light -- the widget's convention.
\newcommand{\qcdark}[2]{\qcglyph{#2}{#1}{qdarkbody}{qdarkedge}}
\newcommand{\qclite}[2]{\qcglyph{#2}{#1}{qlitebody}{qliteedge}}
% \qcgivens{i/j,...} -- reads exactly like the Lean array literal #[(i,j),...].
\newcommand{\qcgivens}[1]{\foreach \qcgi/\qcgj in {#1}{\qcdark{\qcgi}{\qcgj}}}
% \qcplace{q_0,...,q_{N-1}} -- reads exactly like the Lean array #[q_0,...,q_{N-1}].
\newcommand{\qcplace}[1]{%
  \foreach \qccol [count=\qcrow from 0] in {#1}{\qclite{\qcrow}{\qccol}}}
% \qcrule{colour}{i}{j}{i'}{j'} -- one row of A-hat, as a capsule through the centres of its squares.
\newcommand{\qcrule}[5]{%
  \draw[#1,line cap=round,opacity=0.42,line width=0.30cm]
    ({#3+0.5},{-#2-0.5}) -- ({#5+0.5},{-#4-0.5});}
% \qcmark{i}{j} -- outline one square.
\newcommand{\qcmark}[2]{\draw[black,line width=0.9pt] (#2,-#1) rectangle ++(1,-1);}
% \qcattack{i}{j}{i'}{j'} -- the widget's attack indicator: tinted squares and a dashed segment.
\newcommand{\qcattack}[4]{%
  \fill[attackc,opacity=0.28] (#2,-#1) rectangle ++(1,-1);
  \fill[attackc,opacity=0.28] (#4,-#3) rectangle ++(1,-1);
  \draw[attackc,line width=0.8pt,dash pattern=on 2.4pt off 1.4pt]
    ({#2+0.5},{-#1-0.5}) -- ({#4+0.5},{-#3-0.5});}
\newcommand{\qclab}[1]{{\footnotesize #1}}

% Loaded last, and only if the class has not already done so: \secref and \src need \hyperref and
% \href, and an undefined one of those is a hard error rather than a bad box.
\usepackage{hyperref}

\title{Certified Collaborative Neurodynamic Algorithms\\ for Solving $n$-Queens Completion}

\titlerunning{Certified Collaborative Neurodynamic Algorithms for $n$-Queens Completion}

\author{
Matteo Cipollina\inst{1}
\and
Michail Karatarakis\inst{2}
\and
Freek Wiedijk\inst{2}
}

\institute{
  Axiomatic AI, Barcelona, Spain\\
  \email{matteo.cipollina@yahoo.it}
\and
  Radboud University, Nijmegen, The Netherlands\\
  \email{michail.karatarakis@ru.nl, freek@cs.ru.nl}\\
}

\authorrunning{Cipollina, Karatarakis and Wiedijk}

\begin{document}
\maketitle

\begin{abstract}
In this article, $n$-Queens Completion is formulated as a quadratic unconstrained binary
optimization problem, and the formulation is certified in the Lean~4 proof assistant. The objective
is proved to be the energy of a $\{0,1\}$ Boltzmann machine whose minimizers are exactly the
completions of the board, and every infeasible state is proved to lie at least $\tfrac12$ above that
minimum. Soundness and completeness of the formulation are both established, and the Boolean checker
against which they are stated is in turn proved to decide the chess condition written over the
integers. Collaborative neurodynamic optimization algorithms based on discrete Hopfield networks or
Boltzmann machines are then applied unchanged to the formulated problem, a population of such
networks operating concurrently being employed for scatter search and a particle swarm optimization
rule being used to re-initialize their states upon local convergence. Experimental results on eleven
completable and seven blocked instances are elaborated, together with a certified backtracking
baseline, to delimit what the formalization does and does not establish; on six of the seven blocked
instances the absence of a completion is itself a theorem. The variable reduction algorithm of the
Sudoku formulation has no counterpart here, and the differences between the two formulations are
tabulated rather than elided.
\end{abstract}

\section{Introduction}\label{sec:intro}

Discrete Hopfield networks~\citeref{hopfield1982} and Boltzmann machines~\citeref{ackley1985} settle
into states of low energy, and minimizing such an energy is, up to an affine change of variables,
quadratic unconstrained binary optimization (QUBO)~\citeref{lucas2014ising, glover2019}. Constraint
satisfaction problems can therefore be attacked by encoding them as an energy and letting the network
descend it. Collaborative neurodynamic optimization~\citeref{che2019} improves on a single such
network by running a population of them concurrently and re-initializing their states by a particle
swarm rule, and this scheme was applied to Sudoku by Li and Wang~\citeref{liwang2022}.

Two things are missing from that line of work. The encodings are not verified: a network is built
from a problem by hand, and if the construction is wrong then a correct descent returns a wrong
answer, while a descent that finds nothing establishes nothing at all. And the dynamics themselves
are described rather than proved; in particular the momentum recurrences are asserted to converge to
local optima on the strength of simulation.

The first gap is what this article closes, for one problem. In our LPAR-26 paper~\citeref{hnbm} the
two models were formalized in Lean~4 together with their convergence and ergodicity. Here
$n$-Queens Completion is formulated as a QUBO in that framework, the formulation is proved faithful
in both directions, and the collaborative neurodynamic algorithms of~\citeref{liwang2022} are run on
the result without modification.

$n$-Queens Completion is chosen deliberately. Given an $N\times N$ board with some queens already
placed, it asks whether the placement extends to $N$ mutually non-attacking queens. The problem is
\textsf{NP}-complete, and counting its solutions is $\#\textsf{P}$-complete~\citeref{gent2017};
unlike the plain $n$-queens problem, which is solvable for every $N\ge4$ by construction, whether a
completion exists genuinely depends on the givens. A formulation whose negative answers are theorems
is therefore worth having.

\medskip\noindent\textbf{Contributions and overview.} \secref{sec:prelim} recalls the neurodynamic
models and states what is proved about them. \secref{sec:form} formulates the problem, exhibits the
incidence structure of the formulation in \secref{sec:inc}, and tabulates in \secref{sec:diff} every
point at which it differs from the Sudoku formulation of~\citeref{liwang2022}; no correspondence is
claimed elsewhere that is not licensed there. \secref{sec:cert} certifies the formulation and
\secref{sec:stop} marks where the certification stops. \secref{sec:cns} describes the algorithms as
applied, \secref{sec:exp} reports results on the whole eighteen-instance corpus, and \secref{sec:map}
maps every statement of this article to its Lean identifier.

\section{Preliminaries}\label{sec:prelim}

\subsection{Neurodynamic optimization}\label{sec:neuro}

\emph{1) Discrete Hopfield network.} A discrete Hopfield network (DHN) operates in discrete time on
$n$ binary neurons:
\begin{equation}\label{eq:dhn}
  u(t) = Wx(t) - \theta, \qquad x(t+1) = \sigma\bigl(u(t)\bigr),
\end{equation}
where $u\in\mathbb{R}^n$ is the net-input vector, $x\in\{0,1\}^n$ the state vector,
$W\in\mathbb{R}^{n\times n}$ the symmetric connection weight matrix with zero diagonal,
$\theta\in\mathbb{R}^n$ the threshold vector, and $\sigma$ is applied element-wise with
$\sigma(u_i) = 1$ if $u_i > 0$ and $0$ otherwise. The energy
\begin{equation}\label{eq:E}
  E(x) = -\tfrac12\sum_{u\neq v}W_{uv}x_ux_v + \sum_u \theta_u x_u
\end{equation}
never increases when a single neuron is updated, so an asynchronous run reaches a state that no
single update changes.

\emph{2) Momentum variants.} As a variant of the DHN, a DHN with a momentum term (DHNm) is
\begin{equation}\label{eq:dhnm}
  u(t+1) = u(t) + Wx(t) - \theta, \qquad x(t) = \sigma\bigl(u(t)\bigr), \qquad u(0) = 0 .
\end{equation}
A Boltzmann machine replaces the threshold by a stochastic acceptance, and its momentum variant
(BMm) is
\begin{equation}\label{eq:bmm}
  u(t+1) = u(t) + Wx(t) - \theta, \qquad
  \Pr\bigl[x_i(t+1) = 1\bigr] = \frac{1}{1 + e^{-u_i(t)/T(t)}},
\end{equation}
with a geometric cooling schedule $T(t) = T_0\eta^t$ for $0 < \eta < 1$. Equations~\eqref{eq:dhnm}
and~\eqref{eq:bmm} are equations (3) and (6) of~\citeref{liwang2022}, and are the recurrences
actually iterated in \secref{sec:exp}.

\emph{3) Collaborative search.} A population of $N$ such networks is run concurrently. The best state
found is retained as the incumbent $x^{*}$; the initial states are then re-initialized by a particle
swarm rule, and a diversity measure over the population triggers bit-flip mutation when the
population collapses onto $x^{*}$. Termination is $M$ consecutive rounds without improvement to
$x^{*}$.

\subsection{What is proved about these models}\label{sec:hnbm}

Both models of \secref{sec:neuro} are formalized in~\citeref{hnbm} over an abstract network
structure, and that development is used here unchanged: convergence of the asynchronous
single-neuron dynamics, and, for the Boltzmann machine at fixed temperature, existence and
uniqueness of the stationary distribution $\propto e^{-E(x)/T}$, the latter via a formalization of
the Perron--Frobenius theorem.

It must be said at once that this does \emph{not} cover \eqref{eq:dhnm} or \eqref{eq:bmm}. Those
recurrences accumulate their net input, and the accumulated $u$ is not the net input of any state of
the network; the two coincide only at $t = 0$, where $u = 0$. The convergence of DHNm and BMm to
local optima is asserted in~\citeref{liwang2022} on the strength of simulation, and remains
unproved here. What is certified in \secref{sec:cert} is the \emph{objective} these recurrences
descend, not the path they take through the state space; \secref{sec:stop} states this boundary
precisely.

\section{Problem formulation}\label{sec:form}

\subsection{$n$-Queens Completion}\label{sec:qc}

Let $N$ be the board size and let $G \subseteq \{0,\dots,N-1\}^2$ be the set of \emph{givens},
squares already carrying a queen. Let $x_{ij} = 1$ mean that a queen stands on the square in row $i$
and column $j$, and $x_{ij} = 0$ otherwise. A completion is a placement of $N$ queens, no two
attacking, that includes every given. Writing $d$ for a diagonal and $a$ for an anti-diagonal of
length at least two, and $s_d, s_a$ for one auxiliary variable per such diagonal, the formulation is
\begin{subequations}\label{eq:form}
\begin{align}
  \textstyle\sum_{j} x_{ij} &= 1, & i &= 0,\dots,N-1, \label{eq:form-row}\\
  \textstyle\sum_{i} x_{ij} &= 1, & j &= 0,\dots,N-1, \label{eq:form-col}\\
  \textstyle\sum_{(i,j)\in d} x_{ij} + s_d &= 1, & d &= 0,\dots,2N-4, \label{eq:form-diag}\\
  \textstyle\sum_{(i,j)\in a} x_{ij} + s_a &= 1, & a &= 0,\dots,2N-4, \label{eq:form-anti}\\
  x_{ij} &= 1, & (i,j) &\in G, \label{eq:form-giv}\\
  x_{ij}, s_d, s_a &\in \{0,1\}. & & \label{eq:form-bin}
\end{align}
\end{subequations}
Constraints~\eqref{eq:form-row} and~\eqref{eq:form-col} ensure that every row and every column of a
completed board carries exactly one queen. Constraints~\eqref{eq:form-diag}
and~\eqref{eq:form-anti} ensure that no diagonal or anti-diagonal carries two, the auxiliary
variable absorbing the slack when a diagonal is empty. Constraints~\eqref{eq:form-giv} enforce the
givens $G$, and~\eqref{eq:form-bin} the binary requirement.

Only diagonals of length at least two are constrained, a single square being unable to carry two
queens; of the $2N-1$ diagonals in each direction the two extreme ones are single corners, leaving
$2N-3$ in each direction. The number of variables is therefore
\begin{equation}\label{eq:sizes}
  n \;=\; N^2 + 2(2N-3) \;=\; N^2 + 4N - 6 ,
  \qquad
  m \;=\; 6N - 6 + |G| ,
\end{equation}
where $m$ counts the constraints; for $N = 8$ this is $n = 90$ and $m = 42$ with no givens. Both
formulas are confirmed on every row of \tabref{tab:results} and \tabref{tab:blocked}.

\subsection{Quadratic penalties}\label{sec:pen}

Collecting~\eqref{eq:form-row}--\eqref{eq:form-giv} into a $0/1$ matrix
$\hat A \in \{0,1\}^{m\times n}$ and a right-hand side $\hat b = \mathbf{1}$, the quadratic penalty is
\begin{equation}\label{eq:p}
  p(x) \;=\; \lVert \hat Ax - \hat b\rVert_2^2
        \;=\; \sum_{r=1}^{m}\Bigl(\textstyle\sum_{u\in r}x_u - \hat b_r\Bigr)^{2},
\end{equation}
so that $p(x) \ge 0$ always and $p(x) = 0$ exactly when~\eqref{eq:form} is satisfied. Expanding
\eqref{eq:p} and using $x_u^2 = x_u$ to fold the diagonal of $\hat A^{\top}\!\hat A$ into the
threshold gives
\begin{equation}\label{eq:WT}
  W_{uv} = -(\hat A^{\top}\!\hat A)_{uv} \ (u \neq v), \qquad
  W_{uu} = 0, \qquad
  \theta_u = \tfrac12\deg(u) - \textstyle\sum_{r \ni u} \hat b_r ,
\end{equation}
with $\deg(u)$ the number of constraints containing $u$, and then
\begin{equation}\label{eq:Ep}
  E(x) \;=\; \tfrac12\bigl(p(x) - \lVert \hat b\rVert_2^2\bigr) .
\end{equation}
Minimizing the energy~\eqref{eq:E} of the network~\eqref{eq:WT} and solving~\eqref{eq:form} are
therefore the same problem. Since $\hat b \equiv \mathbf{1}$ here, $\theta_u = -\tfrac12\deg(u)$, and
the thresholds are determined by the degrees alone.

\subsection{The incidence structure}\label{sec:inc}

\begin{figure}[t]
\centering
\begin{tabular}{@{}c@{\hspace{2.6em}}c@{}}
\begin{qcpic}{0.62cm}
  \qcboard{4}
  \qcrule{cRow}{1}{0}{1}{3}
  \qcrule{cCol}{0}{2}{3}{2}
  \qcrule{cDiag}{0}{1}{2}{3}
  \qcrule{cAnti}{0}{3}{3}{0}
  \qcmark{1}{2}
  \node[cRow,  font=\scriptsize] at (-0.45,-1.5) {$r_{1}$};
  \node[cDiag, font=\scriptsize] at (1.0,0.34)   {$r_{9}$};
  \node[cCol,  font=\scriptsize] at (2.5,0.34)   {$r_{6}$};
  \node[cAnti, font=\scriptsize] at (3.85,0.34)  {$r_{15}$};
\end{qcpic}
&
\begin{qcpic}{0.62cm}
  \qcboard{4}
  \qcrule{cRow}{0}{0}{0}{3}
  \qcrule{cCol}{0}{0}{3}{0}
  \qcrule{cDiag}{0}{0}{3}{3}
  \qcmark{0}{0}
  \draw[cAnti,densely dotted,line width=0.9pt] (0.5,-0.5) circle (0.36);
  \node[cRow,  font=\scriptsize] at (-0.45,-0.5) {$r_{0}$};
  \node[cCol,  font=\scriptsize] at (0.5,0.34)   {$r_{4}$};
  \node[cDiag, font=\scriptsize] at (3.55,-3.86) {$r_{10}$};
  \node[cAnti, font=\scriptsize] at (1.35,-4.42) {no row};
\end{qcpic}
\\[0.5em]
\qclab{(a) an interior square: $\deg = 4$} & \qclab{(b) a corner: $\deg = 3$}
\end{tabular}
\caption{The rows of $\hat A$ through one square of the $4\times4$ board, on which $n = 26$ and
$m = 18$. Each capsule is one constraint of~\eqref{eq:form}, drawn through the squares whose
variables it contains: \textcolor{cRow}{board row}, \textcolor{cCol}{column},
\textcolor{cDiag}{diagonal}, \textcolor{cAnti}{anti-diagonal}. (a) The square $(1,2)$ lies in
$r_1, r_6, r_9$ and $r_{15}$. (b) The anti-diagonal through the corner $(0,0)$ is the single square
$\{(0,0)\}$, drawn dotted, which generates no constraint, so the corner lies in three rows only. The
indices are those the development assigns.}\label{fig:inc}
\end{figure}

The matrix $\hat A$ of~\eqref{eq:p} is the incidence matrix of a hypergraph on the variables: one row
per constraint, one column per variable, and $\hat A_{ru} = 1$ when variable $u$ occurs in constraint
$r$. Rows are laid out by family --- $N$ board rows, then $N$ columns, then $2N-3$ diagonals, then
$2N-3$ anti-diagonals, then one row per given --- so that board row $i$ is $r_i$ and column $j$ is
$r_{N+j}$. \figref{fig:inc} draws every row through one square at $N = 4$.

The figure also shows why the degrees are not uniform, which is what makes $\theta$ half-integral. An
interior square lies in four rows. A corner lies in three: one of its two diagonals is the corner
itself, and a single square cannot carry two queens, so no constraint is generated for it. An
auxiliary variable lies in exactly one row. Hence $\deg(u) \in \{4,3,1\}$ and, by~\eqref{eq:WT},
$\theta_u \in \{-2,-\tfrac32,-\tfrac12\}$.

One detail of the layout is worth recording, because it is where a transcription error would go
unnoticed. A diagonal is a set of squares on which $i - j$ is constant and an anti-diagonal one on
which $i + j$ is constant; but the development indexes squares by natural numbers, on which
subtraction is truncated, $2 - 5$ being $0$ rather than $-3$. Membership is therefore tested
additively, the square $(i,j)$ lying on the $e$-th row of the diagonal block when $i + N = e + 2 + j$
and on the $e$-th of the anti-diagonal block when $i + j + (2N-3) = e + 1$. Both are equivalent to the
intended condition over $\mathbb{Z}$, and neither ever subtracts. The same hazard appears in the
checker, where it is discharged by \thmref{thm:spec}.

\subsection{How this formulation differs from the Sudoku formulation}\label{sec:diff}

The Sudoku formulation of~\citeref{liwang2022} and the formulation above are both quadratic
unconstrained binary optimizations solved by the same algorithms, but they are not analogous point
for point, and the differences are recorded here so that nothing later depends on a correspondence
that does not hold.

\begin{table}[t]
\centering\small
\setlength{\tabcolsep}{4pt}
\begin{tabular}{@{}p{0.21\textwidth}p{0.34\textwidth}p{0.36\textwidth}@{}}
\toprule
& Sudoku~\citeref{liwang2022} & $n$-Queens Completion (here)\\
\midrule
decision variables & $x_{ijk}$, one per cell and digit; $n^3 = 729$ at $n = 9$ &
  $x_{ij}$, one per square, plus one auxiliary variable per constrained diagonal; $N^2+4N-6$\\
constraint families & four, all \emph{exactly-one}, (10a)--(10d) &
  two exactly-one, two \emph{at-most-one}\\
auxiliary variables & none needed &
  $2(2N-3)$ needed, at-most-one not being an equality\\
column degree of $\hat A$ & uniformly $4$ &
  $4$ at an interior square, $3$ at each of the four corners, $1$ at an auxiliary variable\\
threshold & $\theta = -2\cdot\mathbf{1}$, uniform &
  $\theta_u \in \{-2, -\tfrac32, -\tfrac12\}$, non-uniform and half-integral\\
variable reduction & Algorithm~1; reduces $729$ variables to between $0$ and $28\%$, and solves five
  of ten instances outright & \textbf{no counterpart implemented}; see below\\
givens & enforced by (10e), after reduction &
  enforced as one-variable constraints, with no reduction stage\\
instances reported & five, of ten, after reduction &
  eleven completable boards and seven blocked ones, none reduced\\
negative instances & none; every puzzle posed has a solution &
  seven, on six of which the absence of a solution is a theorem\\
reported statistics & best and worst objective value, mean and standard deviation over runs &
  success counts over seeds; see \secref{sec:stats} for why\\
\bottomrule
\end{tabular}
\caption{Differences between the two formulations.}\label{tab:diff}
\end{table}

\tabref{tab:diff} lists them. Three deserve comment.

\emph{1) At-most-one is not exactly-one.} All four Sudoku families are equalities, so~\eqref{eq:p}
applies directly and no auxiliary variables arise. A diagonal, by contrast, may legitimately be
empty, so \enquote{at most one queen} is not an equality. One auxiliary variable per diagonal
converts it into one: if the diagonal is empty the auxiliary variable is set instead, if it carries
one queen the auxiliary variable is $0$, and if it carries two then no value of the auxiliary
variable satisfies the constraint, so the row contributes at least $1$ to~\eqref{eq:p}. One such
variable suffices per diagonal, whatever the diagonal's length. This is the only structural addition
the queens formulation makes to the Sudoku one, and it is what accounts for the $4N-6$ extra
variables in~\eqref{eq:sizes}.

\emph{2) The degrees are not uniform.} In the Sudoku formulation every variable occurs in exactly
four constraints, so $\theta$ is the constant $-2$. Here the degree is $4$, $3$ or $1$, as
\figref{fig:inc} shows, and $\theta_u$ is correspondingly $-2$, $-\tfrac32$ or $-\tfrac12$.
Half-integral thresholds are harmless for~\eqref{eq:E}, whose $\theta$ is real-valued, but $\theta$
is stored in an integer array so that the inner loop is exact integer arithmetic; what is stored is
therefore $2\theta_u = -\deg(u)$, and it is halved once, where the real-valued energy is defined.

\emph{3) There is no variable reduction algorithm.} Algorithm~1 of~\citeref{liwang2022} is Sudoku
constraint propagation: it fixes the variables of the given cells, eliminates the variables they
exclude in the same row, column and block, and fixes any cell left with a unique candidate. Its
effect is large --- five of the ten instances are solved by it alone. Nothing corresponding is
implemented here, so every result in \secref{sec:exp} is obtained by the collaborative search
operating on the full $N^2+4N-6$ variables. An analogue exists in principle, since a row or column
left with a single square unattacked by a given must carry a queen there, and its absence is a
difference in the pipeline rather than in the formulation. Consequently the queens results are not
comparable with the Sudoku results of~\citeref{liwang2022} in either direction.

\section{Certification of the formulation}\label{sec:cert}

The formulation~\eqref{eq:form}--\eqref{eq:Ep} is developed in Lean~4 on top of the networks
of~\citeref{hnbm}, and the following are proved rather than argued. Throughout, a board is an array
$q$ of $N$ column indices, $q_i$ being the column of the queen in row $i$, so that
\eqref{eq:form-row} holds by construction; this is the representation the figures use.

\subsection{The objective is an energy}\label{sec:energy}

\begin{thm}[the objective is an energy]\label{thm:net}
For any well-formed instance of~\eqref{eq:p}, the choice~\eqref{eq:WT} satisfies~\eqref{eq:Ep};
hence $E$ attains its minimum $-\tfrac12\lVert \hat b\rVert_2^2$ exactly at the points where $p$
vanishes, and, $p$ being integer-valued, every other state lies at least $\tfrac12$ above that
minimum.
\end{thm}

The gap of $\tfrac12$ is what makes the objective a decision procedure rather than an approximation:
no state has energy merely close to the minimum, so a run that ends above the minimum ends at least
$\tfrac12$ above it. It also fixes the reporting convention of \secref{sec:stats}.

\subsection{Soundness and completeness}\label{sec:iff}

\begin{thm}[the formulation is faithful]\label{thm:iff}
Suppose every given lies on the board. If $p(x) = 0$ then the board decoded from $x$ satisfies the
checker; if a board satisfies the checker then the point encoded from it has $p = 0$; and therefore
$p$ has a zero if and only if the instance has a completion.
\end{thm}

The hypothesis is the only one, and in particular no lower bound on $N$ is assumed: at $N = 0$ there
are no variables and no constraints, $p$ is an empty sum, and the empty board is vacuously a
completion. What must be excluded is a given off the board, whose constraint would contain no
variable and could never be satisfied.

\subsection{The checker is the chess condition}\label{sec:spec}

\begin{thm}[the checker is the chess condition]\label{thm:spec}
The Boolean checker of \thmref{thm:iff} decides the condition stated over the integers: that the
board has $N$ entries, all on the board, that no two distinct rows $i \neq k$ satisfy
$q_i = q_k \vee i - q_i = k - q_k \vee i + q_i = k + q_k$, and that every given is present. Hence
$p$ has a zero if and only if a completion exists in that sense.
\end{thm}

\thmref{thm:spec} is not redundant. For the reason given in \secref{sec:inc}, the checker tests the
diagonal conditions in the rearranged form $i + q_k \neq k + q_i$, and it reads the board through an
array accessor with a default value chosen so that a missing entry fails. Neither is a statement
about chess. \thmref{thm:spec} discharges both, and after it the specification a reader has to trust
mentions no arrays, no defaults and no truncated arithmetic. As an additional check the checker is
compared against an independently written implementation over the integers, exhaustively on every
board of size four and of size five.

\subsection{The minimizers are the completions}\label{sec:min}

\begin{thm}[the minimizers are the completions]\label{thm:min}
For the queens instance, a state of the network~\eqref{eq:WT} attains minimum energy if and only if
the board decoded from it is a completion.
\end{thm}

Combining \thmref{thm:net} with the uniqueness result of \secref{sec:hnbm} gives further that the
stationary distribution of the Boltzmann machine at temperature $T$ assigns vanishing probability to
the non-completions as $T \to 0^{+}$. This is a statement about the equilibrium distribution, not
about any cooling schedule.

\subsection{Boards with no completion}\label{sec:blocked}

\begin{figure}[t]
\centering
\begin{tabular}{@{}c@{\hspace{1.7em}}c@{\hspace{1.7em}}c@{\hspace{1.0em}}c@{\hspace{1.0em}}c@{}}
\begin{qcpic}{0.40cm}
  \qcboard{8}\qcattack{0}{0}{1}{1}\qcgivens{0/0,1/1}
\end{qcpic}
&
\begin{qcpic}{0.40cm}
  \qcboard{7}\qcgivens{0/0,1/6}
\end{qcpic}
&
\begin{qcpic}{0.40cm}
  \qcboard{4}\qcgivens{0/0}
\end{qcpic}
&
\begin{qcpic}{0.40cm}
  \qcboard{4}\qcplace{2,0,3,1}
\end{qcpic}
&
\begin{qcpic}{0.40cm}
  \qcboard{4}\qcplace{1,3,0,2}
\end{qcpic}
\\[0.5em]
\qclab{(a) \texttt{Attack}} & \qclab{(b) \texttt{Blk7}} & \qclab{(c) \texttt{Blk4}}
& \qclab{(d)} & \qclab{(e)}
\end{tabular}
\caption{Three blocked instances, and why. Given queens are dark, as in the widget. (a) The two
givens share a diagonal, shown the way the widget shows it, and no enumeration is needed. (b) The
givens of \texttt{Blk7} share no row, no column and no diagonal, and the board admits no completion
regardless; this is the phenomenon that makes Completion hard. (c) A single corner queen at $N = 4$,
refuted by (d) and (e), which are the only two completions of the empty $4\times4$ board: neither
places a queen at $(0,0)$. A corner queen is fatal at $N = 4$ and $N = 6$ and harmless
thereafter.}\label{fig:blocked}
\end{figure}

Because \thmref{thm:iff} is an equivalence, the absence of a zero of $p$ is a statement about the
board. On six of the seven blocked instances of \tabref{tab:blocked} that statement is proved, by
exhausting the tuples the givens leave open using kernel reduction; \lstinline{native_decide} is not
used, so no compiler is part of the trusted base. \figref{fig:blocked} shows three of the six.

\texttt{Blk7} is the instance worth noting. Its two given queens share no row, no column and no
diagonal, so the placement is locally unobjectionable and no attack is displayed, and it admits no
completion. By \thmref{thm:iff} the refutation extends to every one of the $2^{71}$ binary vectors of
that instance, from a check over $7^5 = 16807$ tuples. \texttt{Blk4} illustrates the other way an
instance dies: the empty $4\times4$ board has exactly two completions, drawn as (d) and (e), and
neither uses the corner, so pinning a queen there is fatal. That a corner queen is nonetheless
harmless from $N = 7$ on is visible in \tabref{tab:results}, where \texttt{Comp7} is the same
instance at $N = 7$ and has four completions; the counts for $\langle n, \{(0,0)\}\rangle$ at
$n = 4,\dots,10$ are $0,2,0,4,4,28,64$.

The seventh instance, \texttt{Blk8}, is an $8\times8$ board with two non-attacking givens. It leaves
$8^6 = 262144$ tuples open, beyond what kernel reduction will do here, and it is reported as an
observation of the baseline rather than as a theorem. \tabref{tab:blocked} marks it as such, rather
than presenting it as carrying the same weight as the other six.

\subsection{Where the certification stops}\label{sec:stop}

The objective is identified with the energy by \thmref{thm:net}; the integer routine computing the
net input in the inner loop is proved equal to the network's local field; and the minimizers are
characterized by \thmref{thm:min}. The \emph{trajectory} is not certified, for the reason given in
\secref{sec:hnbm}: \eqref{eq:dhnm} and~\eqref{eq:bmm} accumulate, and no configuration of the
implemented solver is the single-neuron dynamics that~\citeref{hnbm} covers. What a run of
\secref{sec:exp} establishes when it succeeds is therefore a theorem about the board; what it
establishes when it fails is nothing, and the negative results of \secref{sec:blocked} come from
enumeration instead.

\section{Collaborative neurodynamic algorithms}\label{sec:cns}

\subsection{Algorithm 2 as applied}\label{sec:alg}

\begin{figure}[t]
\centering\small
\setlength{\tabcolsep}{0pt}
\begin{tabular}{@{}p{0.965\textwidth}@{}}
\toprule
\textbf{Algorithm 2} (of~\citeref{liwang2022}), as applied to~\eqref{eq:form}.\\
\midrule
\textbf{Input:} the network $(W,\theta)$ of~\eqref{eq:WT}; a model $\mathcal{M} \in
\{\text{DHNm}~\eqref{eq:dhnm},\ \text{BMm}~\eqref{eq:bmm}\}$; and $N$, $M$, $c_0$, $c_1$, $c_2$,
$P_m$, $\mathcal{T}$, $T_0$, $\eta$.\\
\textbf{Output:} the incumbent $x^{*}$ and its objective value $p(x^{*})$.\\
\midrule
\emph{1} \quad draw $x^{(i)}(0)$ with one variable set in each board row, $v^{(i)} \sim
U[-1,1]^{n}$; set $\hat x^{(i)} \leftarrow x^{(i)}(0)$ and $x^{*} \leftarrow \arg\min_i
p(\hat x^{(i)})$\\
\emph{2--7} \quad\, \textbf{repeat:} run each of the $N$ copies of $\mathcal{M}$ from $x^{(i)}(0)$ to
equilibrium, giving $x^{(i)}$; set $\hat x^{(i)} \leftarrow x^{(i)}$ whenever
$p(x^{(i)}) < p(\hat x^{(i)})$\\
\emph{8--14} \quad $x^{*} \leftarrow \arg\min_i p(\hat x^{(i)})$ if that improves on $x^{*}$, and
reset the stall counter; otherwise increment it. Exit if $p(x^{*}) = 0$\\
\emph{15--20} \, $v^{(i)} \leftarrow c_0 v^{(i)} + c_1 r_1 (\hat x^{(i)} - x^{(i)})
+ c_2 r_2 (x^{*} - x^{(i)})$ with $r_1, r_2 \sim U[0,1]$, then
$x^{(i)}(0) \leftarrow \mathrm{round}\bigl(x^{(i)}(0) + v^{(i)}\bigr)$\\
\emph{21--24} \, $\delta \leftarrow \frac{1}{Nn}\sum_i \lVert \hat x^{(i)} - x^{*}\rVert_2$;
\ \textbf{if} $\delta\sqrt{n} < \mathcal{T}$ \textbf{then} flip each bit of every $x^{(i)}(0)$ with
probability $P_m$\\
\quad\ \ \ \ \ \ \ \, \textbf{until} $p(x^{*}) = 0$, or $M$ consecutive rounds pass without
improvement, or the round cap is reached\\
\bottomrule
\end{tabular}
\caption{The collaborative search, with the step numbers of~\citeref{liwang2022}. The population is
re-initialized from the \emph{equilibrium} states $x^{(i)}$ at steps 15--20 while the personal bests
$\hat x^{(i)}$ supply the cognitive term, and the diversity $\delta$ of step 21 is measured over the
personal bests. The one-hot draw at step 1, the $\sqrt{n}$ at step 22 and the target of the mutation at
step 23 are the three points a literal reading of~\citeref{liwang2022} does not settle; all three are
discussed in \secref{sec:params}.}\label{fig:alg}
\end{figure}

Algorithm~2 of~\citeref{liwang2022}, reproduced as \figref{fig:alg}, is applied without modification
to its structure. A population of $N$ copies of DHNm \eqref{eq:dhnm} or BMm \eqref{eq:bmm} is run to
equilibrium; the incumbent $x^{*}$ is updated from the best state found; the initial states are
re-initialized by the particle swarm rule; and bit-flip mutation fires when the diversity of the
population falls below a threshold $\mathcal{T}$. The only problem-specific input is the list of
mutually exclusive variable groups used at step 1, which for this formulation is the $N$ squares of
each board row, that is to say the rows $r_0,\dots,r_{N-1}$ of \figref{fig:inc}. As recorded in
\secref{sec:diff}, no variable reduction precedes the search.

\subsection{Parameters recovered}\label{sec:params}

\begin{table}[t]
\centering\small
\setlength{\tabcolsep}{5pt}
\begin{tabular}{@{}llll@{}}
\toprule
& meaning & value & \\
\midrule
$N$            & population size                            & $20$ & \\
$M$            & stall limit, in rounds                     & $20$ & \\
               & round cap                                  & $60$ & \\
               & inner iterations per equilibration         & $40$ & \\
$T_0$          & initial temperature                        & $3.0$ & \\
$\eta$         & cooling rate                               & $0.9$ & \\
$c_0$          & inertia weight                             & $0.5$ & \\
$c_1$          & cognitive coefficient                      & $2.0$ & \\
$c_2$          & social coefficient                         & $0.25$ & decisive; see below\\
$P_m$          & bit-flip mutation probability              & $0.05$ & \\
$\mathcal{T}$  & diversity threshold, as a fraction of $\delta_{\max} = 1/\sqrt{n}$
               & $0.4$ (BMm), $0.9$ (DHNm) & scale-free; see below\\
\bottomrule
\end{tabular}
\caption{The parameters of \figref{fig:alg}, as used for every row of \tabref{tab:results} and
\tabref{tab:blocked}. \emph{Not one of these values is stated in~\citeref{liwang2022}}; all were
recovered experimentally, and the two annotated ones are not merely a matter of
tuning.}\label{tab:params}
\end{table}

Reproducing~\citeref{liwang2022} requires numerical values for every parameter in
\tabref{tab:params}. None is stated there. Two of them are not merely a matter of tuning.

\emph{1) The social coefficient.} The coefficient $c_2$ weights the pull of the incumbent. If it is
comparable to $c_1$ then every copy is drawn onto $x^{*}$, and once $x^{*}$ lies on a plateau --- which
on these instances happens almost at once --- the population ceases to explore, and further copies are
further copies of the same stuck point. Taking $c_2$ well below $c_1$ removes this.

\emph{2) The diversity threshold.} $\mathcal{T}$ must be made scale-free. The diversity measure
of~\citeref{liwang2022} is an average of Euclidean distances divided by the number of variables, and
is therefore at most $1/\sqrt{n}$; a fixed threshold consequently means something different on every
instance, and on a large one fires on every round. The threshold is compared here against the
diversity scaled by $\sqrt{n}$, as step 22 of \figref{fig:alg} records. The recovered value is
model-specific, $0.9$ for DHNm against $0.4$ for BMm; having been fitted on Sudoku it transfers, in
that on the empty $8\times8$ board DHNm fails at BMm's $0.4$ and succeeds on nine of ten seeds at
$0.9$.

Two further choices are interpretations rather than numbers, and are recorded because
\citeref{liwang2022} does not settle them. Step 1 there says only that
$x^{(i)}(0) \in \{0,1\}^{n^2}$; uniform bits set about half the variables, whereas a feasible point
sets exactly one per board row, so a uniformly seeded run spends its whole equilibration merely
reaching one queen per row. Seeding one-hot starts every copy on the feasible set
of~\eqref{eq:form-row}, leaving the dynamics the columns and the diagonals. And step 23 says only to
mutate \enquote{according to Eq. (9)} without naming the target; the mutation is applied here to the
seeds $x^{(i)}(0)$, and never to $x^{*}$.

\section{Experimental results}\label{sec:exp}

\subsection{Setup}\label{sec:setup}

The parameters of \tabref{tab:params} are fixed across every instance, and the base seed is fixed too.
The generator is \texttt{splitmix64} and no clock is read, so both tables are reproducible exactly.
The baseline is a chronological backtracking completion search written independently of the
formulation. It is itself certified, in that every board it returns is proved to satisfy the checker,
and the column \textsf{ok} of \tabref{tab:results} is that theorem being re-checked at run time. The
converse is not proved, so the baseline's failures are reported as observations.

\subsection{Results}\label{sec:results}

\begin{table}[t]
\centering\footnotesize
\setlength{\tabcolsep}{4.5pt}
\begin{tabular}{@{}lrrrr rr rrr r r@{}}
\toprule
& & & & & \multicolumn{2}{c}{backtracking}
& \multicolumn{3}{c}{CNS/BMm, \eqref{eq:bmm}\,+\,Alg.~2} & CNS/DHNm & BMm alone\\
\cmidrule(lr){6-7}\cmidrule(lr){8-10}\cmidrule(lr){11-11}\cmidrule(lr){12-12}
board & $N$ & $|G|$ & $n$ & $m$ & ok & \#compl. & hit & rounds & $s/k$ & $s/k$ & $s/20$\\
\midrule
\texttt{Q4}     & 4  & 0 & 26  & 18 & y & 2   & y & 1  & 10/10 & 10/10 & 7\\
\texttt{Q5}     & 5  & 0 & 39  & 24 & y & 10  & y & 1  & 10/10 & 10/10 & 5\\
\texttt{Q6}     & 6  & 0 & 54  & 30 & y & 4   & y & 4  & 9/10  & 8/10  & 0\\
\texttt{Q7}     & 7  & 0 & 71  & 36 & y & 40  & y & 10 & 10/10 & 10/10 & 0\\
\texttt{Q8}     & 8  & 0 & 90  & 42 & y & 92  & y & 3  & 10/10 & 9/10  & 0\\
\texttt{Q9}     & 9  & 0 & 111 & 48 & y & 352 & y & 24 & 4/5   & 5/5   & 0\\
\texttt{Q10}    & 10 & 0 & 134 & 54 & y & 724 & n & 31 & 0/5   & 1/5   & 0\\
\midrule
\texttt{Comp6}  & 6  & 1 & 54  & 31 & y & 1   & n & 22 & 6/10  & 7/10  & 0\\
\texttt{Comp7}  & 7  & 1 & 71  & 37 & y & 4   & y & 8  & 10/10 & 6/10  & 0\\
\texttt{Comp8}  & 8  & 3 & 90  & 45 & y & 1   & y & 14 & 6/10  & 3/10  & 0\\
\texttt{Comp8b} & 8  & 5 & 90  & 47 & y & 1   & y & 4  & 10/10 & 10/10 & 0\\
\bottomrule
\end{tabular}
\caption{The eleven completable instances. Columns: board size; number of givens; the variable and
constraint counts $n$ and $m$ of~\eqref{eq:sizes}; whether the checker accepts the board
backtracking returned; the total number of completions; whether CNS/BMm reached $p = 0$ at the base
seed, and in how many rounds; successes over $k$ seeds for CNS/BMm and for CNS/DHNm; and successes
over $20$ seeds for a single BMm run with no population. Backtracking found a completion on every
row.}\label{tab:results}
\end{table}

\begin{table}[t]
\centering\footnotesize
\setlength{\tabcolsep}{4.5pt}
\begin{tabular}{@{}lrrrr rr rr@{}}
\toprule
board & $N$ & $|G|$ & $n$ & $m$ & \#compl. & refuted & tuples & best $p$\\
\midrule
\texttt{Q2}     & 2 & 0 & 6  & 6  & 0 & y & $2^2 = 4$      & 1\\
\texttt{Q3}     & 3 & 0 & 15 & 12 & 0 & y & $3^3 = 27$     & 1\\
\texttt{Blk4}   & 4 & 1 & 26 & 19 & 0 & y & $4^3 = 64$     & 1\\
\texttt{Blk6}   & 6 & 1 & 54 & 31 & 0 & y & $6^5 = 7776$   & 1\\
\texttt{Blk7}   & 7 & 2 & 71 & 38 & 0 & y & $7^5 = 16807$  & 1\\
\texttt{Attack} & 8 & 2 & 90 & 44 & 0 & y & none needed    & 1\\
\texttt{Blk8}   & 8 & 2 & 90 & 44 & 0 & n & $8^6 = 262144$ & 1\\
\bottomrule
\end{tabular}
\caption{The seven blocked instances. \textsf{refuted} is whether the absence of a completion is a
theorem of the development rather than a verdict of the baseline; \textsf{tuples} is the enumeration
that establishes it, over the rows the givens leave open. \texttt{Attack} needs none, its two givens
sharing a diagonal. On every row the search found no zero on any seed, and its best value was exactly
$1$, the least it can be by \thmref{thm:net}.}\label{tab:blocked}
\end{table}

\tabref{tab:results} and \tabref{tab:blocked} report the whole corpus, unreduced and untruncated.
Four observations.

\emph{1) The baseline wins.} Backtracking solves every completable instance and refutes every blocked
one, in times not worth tabulating. No speed claim is made for the neurodynamic algorithms, and
nothing below should be read as one.

\emph{2) The population is not a refinement.} A single BMm run with no population succeeds on $7$ of
$20$ seeds at $N = 4$ and $5$ of $20$ at $N = 5$, and on nothing from $N = 6$ upwards or on any
instance with givens. The collaborative search reaches a completion on ten of the eleven rows. Here
the collaborative element is the difference between working and not working.

\emph{3) The two models separate, and only where there are givens.} On the seven empty boards DHNm
and BMm are indistinguishable, at $53$ of $60$ seeds each. On the four boards with givens BMm is
ahead, $32$ of $40$ against $26$. The separation is not uniform --- DHNm is the only configuration to
reach $N = 10$ at all, once in five seeds --- and it is problem-specific: on a graph-colouring corpus
formulated in the same framework the two are indistinguishable throughout.

\emph{4) The corpus is what bounds the claim.} $N = 10$ is where the swarm stops. A purpose-built
Hopfield network for $n$-queens~\citeref{mandziuk1995} reaches $N = 80$, on the plain problem rather
than on Completion. The completion counts in \tabref{tab:results} are a check on the baseline, which
returns $1,1,0,0,2,10,4,40,92,352,724$ for $n = 0,\dots,10$ on the empty boards, the classical
sequence.

\subsection{Why success counts and not means}\label{sec:stats}

\citeref{liwang2022} reports best and worst objective values with means and standard deviations. That
statistic is uninformative for this formulation. By \thmref{thm:net} the objective is integer-valued,
with a gap of $\tfrac12$ in the energy, so a run ends either at $p = 0$ or at $p \ge 1$; and on every
blocked instance every failed run ends at exactly $1$, as the last column of \tabref{tab:blocked}
records. A mean over runs therefore reports little beyond the success rate, which is given directly,
and a standard deviation reports the same thing twice.

\subsection{Visualization}\label{sec:vis}

\begin{figure}[t]
\centering
\begin{tabular}{@{}c@{\hspace{1.2em}}c@{\hspace{2.6em}}c@{\hspace{1.2em}}c@{}}
\begin{qcpic}{0.40cm}
  \qcboard{6}\qcgivens{0/1}
\end{qcpic}
&
\begin{qcpic}{0.40cm}
  \qcboard{6}\qcplace{1,3,5,0,2,4}\qcgivens{0/1}
\end{qcpic}
&
\begin{qcpic}{0.40cm}
  \qcboard{8}\qcgivens{0/0,1/4,2/7}
\end{qcpic}
&
\begin{qcpic}{0.40cm}
  \qcboard{8}\qcplace{0,4,7,5,2,6,1,3}\qcgivens{0/0,1/4,2/7}
\end{qcpic}
\\[0.5em]
\qclab{(a) \texttt{Comp6}} & \qclab{(b) its completion}
& \qclab{(c) \texttt{Comp8}} & \qclab{(d) its completion}
\end{tabular}
\caption{Two completable instances and their completions, in the colours the widget uses: given
queens dark, queens the solver placed light. Each instance has exactly one completion, so what is
drawn in (b) and (d) is the unique board on which $p$ vanishes. (d) is also the lexicographically
first of the $92$ solutions of the empty $8\times8$ board, pinned by its three
givens.}\label{fig:comp}
\end{figure}

A run may be watched in the Lean infoview as a chessboard, one frame per round, drawn exactly as in
\figref{fig:blocked} and \figref{fig:comp}: given queens dark, placed queens light, and a dashed
segment joining every attacking pair as in \figref{fig:blocked}(a). A segment appears exactly when one
of the pairwise conditions of \thmref{thm:spec} fails, which is exactly when a column or diagonal
constraint of $\hat A$ carries two set variables, so what is watched disappearing is the residual
of~\eqref{eq:p}. One widget is interactive: queens are placed by clicking and the solver is then run
on the resulting instance through the proof assistant's remote-procedure-call mechanism, so the code
that runs is the code the theorems are about.

\section{From the formulation to the development}\label{sec:map}

\begin{table}[t]
\centering\footnotesize
\setlength{\tabcolsep}{4pt}
\begin{tabular}{@{}p{0.39\textwidth}p{0.30\textwidth}l@{}}
\toprule
statement & Lean identifier & file\\
\midrule
\eqref{eq:p}, the canonical form & \lstinline{Problem} & \texttt{QUBO/Problem.lean}\src{HopfieldNet/QUBO/Problem.lean}\\
$p$ is integer-valued & \lstinline{one_le_penaltyR_of_ne_zero} & \texttt{QUBO/Minimizers.lean}\src{HopfieldNet/QUBO/Minimizers.lean}\\
\thmref{thm:net}, \eqref{eq:Ep} & \lstinline{zeroOneHamiltonian_eq} & \texttt{QUBO/Net.lean}\src{HopfieldNet/QUBO/Net.lean}\\
minimizers are the feasible points & \lstinline{energy_eq_min_iff} & \texttt{QUBO/Minimizers.lean}\\
the gap of $\tfrac12$ & \lstinline{half_le_energy_sub_min} & \texttt{QUBO/Minimizers.lean}\\
the inner loop is the local field & \lstinline{netVec_eq_localField} & \texttt{QUBO/Refine.lean}\src{HopfieldNet/QUBO/Refine.lean}\\
low-temperature concentration & \lstinline{stationary_infeasible_tendsto_zero} & \texttt{QUBO/Gibbs.lean}\src{HopfieldNet/QUBO/Gibbs.lean}\\
\midrule
\eqref{eq:form} is well formed & \lstinline{problem_wf}, \lstinline{problem_refines} & \texttt{Queens/Basic.lean}\src{HopfieldNet/QUBO/Instances/Queens/Basic.lean}\\
\figref{fig:inc}, the row layout & \lstinline{rowRow}, \lstinline{colRow}, \lstinline{diagRow}, \lstinline{antiRow} & \texttt{Queens/Basic.lean}\\
additive diagonal membership & \lstinline{onSlack} & \texttt{Queens/Basic.lean}\\
\figref{fig:inc}, the degrees $4$, $3$, $1$ & \lstinline{rowsOf}, and the \lstinline{#guard}s pinning its size & \texttt{Queens/Basic.lean}, \texttt{Spec.lean}\src{HopfieldNet/QUBO/Instances/Queens/Spec.lean}\\
\thmref{thm:iff}, soundness & \lstinline{decode_isQueens} & \texttt{Queens/Basic.lean}\\
\thmref{thm:iff}, completeness & \lstinline{encode_penalty_zero} & \texttt{Queens/Basic.lean}\\
\thmref{thm:iff}, the equivalence & \lstinline{exists_zero_iff_queens} & \texttt{Queens/Basic.lean}\\
\thmref{thm:spec} & \lstinline{isQueens_iff}, \lstinline{exists_zero_iff_completion} & \texttt{Queens/Spec.lean}\\
agreement with an independent checker & \lstinline{isQueens_eq_naive_four/_five} & \texttt{Queens/Spec.lean}\\
\thmref{thm:min} & \lstinline{exists_minEnergy_iff} & \texttt{Queens/Converge.lean}\src{HopfieldNet/QUBO/Instances/Queens/Converge.lean}\\
\figref{fig:blocked}(a), (c) & \lstinline{attacking_no_queens}, \lstinline{blocked4_no_queens} & \texttt{Queens/Examples.lean}\src{HopfieldNet/QUBO/Instances/Queens/Examples.lean}\\
\figref{fig:blocked}(b) & \lstinline{blk7_no_queens}, \lstinline{blk7_no_zero} & \texttt{Queens/Bench.lean}\src{HopfieldNet/QUBO/Instances/Queens/Bench.lean}\\
\midrule
the baseline returns only completions & \lstinline{solve_isQueens} & \texttt{Queens/Baseline.lean}\src{HopfieldNet/QUBO/Instances/Queens/Baseline.lean}\\
\figref{fig:blocked}(d), (e), and the counts & \lstinline{count} & \texttt{Queens/Baseline.lean}\\
Algorithm~2 & \lstinline{search} & \texttt{QUBO/Search.lean}\src{HopfieldNet/QUBO/Search.lean}\\
\eqref{eq:dhnm}, \eqref{eq:bmm} & \lstinline{dhnmStep}, \lstinline{bmmStep} & \texttt{QUBO/Dynamics.lean}\src{HopfieldNet/QUBO/Dynamics.lean}\\
\tabref{tab:results}, \tabref{tab:blocked} & \lstinline{table}, \lstinline{tableBlocked} & \texttt{Queens/Bench.lean}\\
\figref{fig:comp}, the two instances & \lstinline{small6}, \lstinline{completion8} & \texttt{Queens/Examples.lean}\\
every board drawn in this article & the \lstinline{#guard}s of \secref{sec:map} & \texttt{Queens/Figures.lean}\src{HopfieldNet/QUBO/Instances/Queens/Figures.lean}\\
the widgets & \lstinline{animate}, \lstinline{board} & \texttt{Queens/Widget.lean}, \texttt{Interactive.lean}\src{HopfieldNet/QUBO/Instances/Queens/Interactive.lean}\\
\bottomrule
\end{tabular}
\caption{Each statement of this article, its Lean identifier, and its file.}\label{tab:map}
\end{table}

\tabref{tab:map} lists every statement made above against the identifier that discharges it.

The figures are included, and in their case the correspondence is checked rather than asserted. A
chessboard in a paper is a claim about the development that no proof assistant sees, so every board
drawn above is restated in the development as a \lstinline{#guard}, which fails during elaboration if
it stops holding: that \figref{fig:blocked}(b) is the instance \texttt{Blk7} with those two givens,
that they do not attack each other, and that the board admits no completion; that
\figref{fig:blocked}(d) and (e) are the only two completions of the empty $4\times4$ board and that
neither uses the corner; that the boards of \figref{fig:comp} are the unique completions of the
instances beside them; and that each capsule of \figref{fig:inc} covers exactly the squares of the row
of $\hat A$ it is labelled with. The size formulas~\eqref{eq:sizes} are checked there against both
corpora at once rather than read off the printed tables. If a figure and the code ever disagree, the
build stops.

\section{Concluding remarks}\label{sec:concl}

In this article, $n$-Queens Completion is formulated as a quadratic unconstrained binary
optimization problem, and that formulation is certified in Lean~4: it is proved sound and complete,
the checker against which those properties are stated is proved to decide the chess condition
written over the integers, and the objective is proved to be the energy of a $\{0,1\}$ Boltzmann
machine whose minimizers are exactly the completions. Collaborative neurodynamic optimization
algorithms based on a population of discrete Hopfield networks or Boltzmann machines are applied to
the formulated problem without modification, and experimental results on eleven completable and
seven blocked instances are elaborated, six of the latter refuted by proof. No variable reduction
algorithm is proposed, and the differences between this formulation and the Sudoku formulation
of~\citeref{liwang2022} are tabulated in \tabref{tab:diff} rather than assumed away.

What the certification buys is not speed. A backtracking search outperforms the neurodynamic
algorithms on every instance, and a purpose-built Hopfield network for
$n$-queens~\citeref{mandziuk1995} outperforms them at sizes this formulation does not reach. What it
buys is that a run which ends without a zero has established a property of the chessboard, and that
the objective being descended is provably the energy of the network said to be descending it.

Two directions remain. The trajectory is not certified, as \secref{sec:stop} sets out:
\eqref{eq:dhnm} and~\eqref{eq:bmm} accumulate their net input and are not the single-neuron dynamics
for which~\citeref{hnbm} proves convergence, so a solver restricted to that dynamics would be
certified from the board to the last update, and none is. And the variable reduction stage absent
from \tabref{tab:diff} would, if implemented and certified, make the comparison
with~\citeref{liwang2022} a fair one; at present it is not.

\section*{Acknowledgements}
% TODO

\begin{thebibliography}{99}

\bibitem{hnbm}
M.~Cipollina, M.~Karatarakis, and F.~Wiedijk.
\newblock Formalized Hopfield networks and Boltzmann machines.
\newblock In \emph{LPAR-26}, EPiC Series in Computing, 2026. To appear. arXiv:2512.07766.

\bibitem{liwang2022}
H.~Li and J.~Wang.
\newblock Collaborative neurodynamic algorithms for solving Sudoku puzzles.
\newblock In \emph{12th International Conference on Information Science and Technology (ICIST)},
pages 8--17. IEEE, 2022.

\bibitem{che2019}
H.~Che and J.~Wang.
\newblock A collaborative neurodynamic approach to global and combinatorial optimization.
\newblock \emph{Neural Networks}, 114:15--27, 2019. DOI 10.1016/j.neunet.2019.02.002.

\bibitem{hopfield1982}
J.~J.~Hopfield.
\newblock Neural networks and physical systems with emergent collective computational abilities.
\newblock \emph{Proceedings of the National Academy of Sciences}, 79(8):2554--2558, 1982.

\bibitem{ackley1985}
D.~H.~Ackley, G.~E.~Hinton, and T.~J.~Sejnowski.
\newblock A learning algorithm for Boltzmann machines.
\newblock \emph{Cognitive Science}, 9(1):147--169, 1985.

\bibitem{gent2017}
I.~P.~Gent, C.~Jefferson, and P.~Nightingale.
\newblock Complexity of $n$-Queens Completion.
\newblock \emph{Journal of Artificial Intelligence Research}, 59:815--848, 2017.

\bibitem{mandziuk1995}
J.~Ma\'{n}dziuk.
\newblock Solving the $N$-Queens problem with a binary Hopfield-type network.
\newblock \emph{Biological Cybernetics}, 72:439--445, 1995. DOI 10.1007/BF00201419.

\bibitem{lucas2014ising}
A.~Lucas.
\newblock Ising formulations of many NP problems.
\newblock \emph{Frontiers in Physics}, 2:5, 2014. arXiv:1302.5843.

\bibitem{glover2019}
F.~Glover, G.~Kochenberger, and Y.~Du.
\newblock Quantum bridge analytics I: a tutorial on formulating and using QUBO models.
\newblock arXiv:1811.11538, 2019.

\end{thebibliography}

\end{document}
